% ============================================================
%  فصل ۸: مجلس خبرگان قانون اساسی در برابر پیش‌نویس حبیبی-کاتوزیان
% ============================================================
\chapter{مجلس خبرگان قانون اساسی در برابر پیش‌نویس حبیبی-کاتوزیان}

\begin{infoBox}[چکیدهٔ فصل]
فرایند تدوین قانون اساسی جمهوری اسلامی، یکی از \keyword{مهم‌ترین بزنگاه‌ها} در بحث «وعده یا خدعه» است. پیش‌نویسی که تیمی از حقوقدانان (شامل دکتر حسن حبیبی، دکتر هما کاتوزیان، و دکتر ناصر کاتوزیان) تهیه کرده بودند، یک \keyword{جمهوری پارلمانی دموکراتیک} بدون ولایت فقیه بود — کاملاً سازگار با وعده‌های پاریس. اما این پیش‌نویس کنار گذاشته شد و به‌جای آن، مجلس خبرگان قانون اساسی تشکیل شد — مجلسی با اکثریت روحانی که ولایت فقیه را به محور قانون اساسی تبدیل کرد. این فصل به بازسازی تفصیلی این فرایند، بازیگران آن، و تفسیرهای رقیب از آن می‌پردازد.
\end{infoBox}

% ────────────────────────────────────────
\section{پیش‌نویس اولیه: «قانون اساسی بدون ولایت فقیه»}
% ────────────────────────────────────────

\subsection{تیم تدوین‌کننده}

پس از پیروزی انقلاب، بازرگان (نخست‌وزیر دولت موقت) از تیمی از حقوقدانان خواست پیش‌نویس قانون اساسی را تهیه کنند. اعضای اصلی این تیم عبارت بودند از:

\begin{longtable}{%
    >{\raggedleft\arraybackslash}p{0.18\textwidth}
    >{\raggedleft\arraybackslash}p{0.25\textwidth}
    >{\raggedleft\arraybackslash}p{0.25\textwidth}
    >{\raggedleft\arraybackslash}p{0.22\textwidth}}
\toprule
\rowcolor{cPrimary!10}
\textbf{شخص} & \textbf{تخصص} & \textbf{سابقه} & \textbf{موضع سیاسی} \\
\midrule
\endhead

دکتر حسن حبیبی &
دکترای حقوق از فرانسه &
استاد دانشگاه، حقوقدان &
ملی-مذهبی، میانه‌رو \\
\midrule

دکتر هما کاتوزیان &
دکترای اقتصاد و تاریخ از \en{Oxford} &
استاد دانشگاه، مورخ &
لیبرال-دموکرات \\
\midrule

دکتر ناصر کاتوزیان &
دکترای حقوق مدنی &
استاد حقوق دانشگاه تهران &
حقوقدان مستقل \\
\midrule

احمد صدر حاج سید جوادی &
حقوقدان و روزنامه‌نگار &
عضو جبههٔ ملی &
ملی-دموکرات \\

\bottomrule
\caption{تیم تدوین‌کنندهٔ پیش‌نویس قانون اساسی}
\label{tab:draft-team}
\end{longtable}

\subsection{ویژگی‌های پیش‌نویس}

پیش‌نویس حبیبی-کاتوزیان در خرداد ۱۳۵۸ آماده شد و در ۲۶ خرداد ۱۳۵۸ توسط دولت موقت منتشر گردید.\footnote{\en{Schirazi (1997), pp.\,22--26;} همچنین: \en{Arjomand (1988), pp.\,150--155.} متن کامل پیش‌نویس در روزنامهٔ «اطلاعات» و «کیهان» خرداد ۱۳۵۸ منتشر شد.} ویژگی‌های اصلی:

\begin{principleBox}[ارکان پیش‌نویس حبیبی-کاتوزیان]
\begin{enumerate}[label=\textcolor{cPrimary}{\bfseries\arabic*.}, rightmargin=2em]
    \item \textbf{رئیس‌جمهور منتخب مردم} — بالاترین مقام کشور.
    \item \textbf{مجلس شورای ملی} — قوهٔ مقننه با اختیارات کامل.
    \item \textbf{قوهٔ قضاییه مستقل} — با قضات حرفه‌ای.
    \item \textbf{شورای نگهبان محدود} — فقط بررسی انطباق قوانین با موازین شرعی (مشابه اصل دوم متمم قانون اساسی مشروطه).
    \item \textbf{بدون ولایت فقیه} — هیچ اشاره‌ای به ولایت یا نظارت فقیه بر کل نظام.
    \item \textbf{آزادی احزاب و مطبوعات} — بدون قید «موازین اسلامی».
    \item \textbf{حقوق اقلیت‌ها} — برابری کامل شهروندی.
    \item \textbf{ارتش ملی} — تحت فرماندهی رئیس‌جمهور، بدون سپاه پاسداران.
\end{enumerate}
\end{principleBox}

\begin{noteBox}[سازگاری پیش‌نویس با وعده‌های پاریس]
پیش‌نویس حبیبی-کاتوزیان دقیقاً همان نظامی بود که خمینی در پاریس وعده‌اش را داده بود: جمهوری دموکراتیک، حکومت متخصصان، بدون حکومت روحانیون، با آزادی‌های مدنی. اگر خمینی واقعاً به وعده‌هایش پایبند بود، باید از این پیش‌نویس حمایت می‌کرد. آنچه رخ داد، خلاف این بود.
\end{noteBox}

\subsection{واکنش اولیهٔ خمینی: تأیید مشروط}

در خرداد ۱۳۵۸، خمینی ابتدا واکنش مثبتی به پیش‌نویس نشان داد. بازرگان نقل می‌کند:

\begin{warningBox}[روایت بازرگان]
\textit{«پیش‌نویس را خدمت آقا [خمینی] بردیم. ایشان مطالعه فرمودند و فرمودند: "خوب است. چرا به رفراندوم نمی‌گذارید؟"»}\footnote{مهدی بازرگان، \textit{انقلاب ایران در دو حرکت}، ص ۱۲۵–۱۲۸. همچنین روایت مشابهی در خاطرات یزدی آمده.}
\end{warningBox}

اما ظرف چند هفته، موضع خمینی تغییر کرد. چه اتفاقی افتاد؟

% ────────────────────────────────────────
\section{چرخش: از رفراندوم مستقیم به مجلس خبرگان}
% ────────────────────────────────────────

\subsection{فشار حزب جمهوری اسلامی}

در فاصلهٔ خرداد تا مرداد ۱۳۵۸، حزب جمهوری اسلامی — به رهبری بهشتی، باهنر، و خامنه‌ای — کارزاری علیه رفراندوم مستقیم بر پیش‌نویس به راه انداخت. استدلال‌های آن‌ها:

\begin{itemize}[label=\textcolor{cAccent}{$\blacktriangleright$}, rightmargin=2em]
    \item «پیش‌نویس بدون بحث کارشناسی به رفراندوم گذاشته شود عجولانه است.»
    \item «نخبگان حوزوی باید بررسی کنند آیا با اسلام سازگار است.»
    \item «مردم حق دارند از طریق نمایندگان‌شان دربارهٔ قانون اساسی بحث کنند.»
    \item «مجلس مؤسسان (خبرگان) باید تشکیل شود.»\footnote{\en{Schirazi (1997), pp.\,26--30.} همچنین: \en{Martin (2000), pp.\,140--148.}}
\end{itemize}

\begin{warningBox}[تحلیل انتقادی]
استدلال‌های حزب جمهوری اسلامی ظاهراً «دموکراتیک» بودند (بحث کارشناسی، نمایندگی مردم). اما نتیجه خلاف دموکراسی بود:

\begin{enumerate}[label=\textcolor{cAccent}{\bfseries\arabic*.}, rightmargin=2em]
    \item مجلس خبرگان ۷۳ نفره بود — در حالی‌که پیشنهاد اولیهٔ بازرگان «مجلس مؤسسان ۳۰۰ نفره» بود. عدد کوچک‌تر = کنترل آسان‌تر.\footnote{بازرگان خواستار مجلس مؤسسان ۳۰۰ نفره بود. خمینی ابتدا موافقت کرد اما سپس عدد به ۷۳ تقلیل یافت — ظاهراً بنا بر پیشنهاد بهشتی. \en{Bakhash (1984), pp.\,72--75.}}
    \item انتخابات مجلس خبرگان (۱۲ مرداد ۱۳۵۸) در شرایطی برگزار شد که:
    \begin{itemize}
        \item بسیاری از نامزدهای غیرروحانی مورد حمله قرار گرفتند.
        \item حزب جمهوری اسلامی تنها تشکیلات سازمان‌یافتهٔ سراسری بود.
        \item میزان مشارکت نسبتاً پایین بود (حدود ۵۱\%).\footnote{\en{Arjomand (1988), pp.\,155--158.}}
    \end{itemize}
    \item نتیجه: ۵۵ نفر از ۷۳ نفر روحانی بودند و اکثریت از حزب جمهوری اسلامی.
\end{enumerate}
\end{warningBox}

\subsection{نقش خمینی در چرخش}

پرسش محوری: آیا خمینی خودش تصمیم به تغییر مسیر گرفت، یا تحت فشار حلقهٔ درونی بود؟

\begin{principleBox}[چهار خوانش از «چرخش»]
\begin{description}[style=nextline, font=\bfseries\color{cPrimary}]
    \item[تز ۱ (صداقت + اجبار)]
    خمینی ابتدا واقعاً موافق رفراندوم بود. اما وقتی دید گروه‌های مارکسیست و لیبرال‌های سکولار دارند سازمان‌دهی می‌کنند، نگران شد و پذیرفت مجلس خبرگان تشکیل شود تا «اسلامیت» نظام حفظ شود.

    \item[تز ۲ (خدعه)]
    خمینی از ابتدا قصد رفراندوم مستقیم نداشت. تأیید اولیهٔ پیش‌نویس «ژست» بود. وقتی حزب جمهوری اسلامی بهانهٔ لازم را فراهم کرد، خمینی مسیر واقعی‌اش را دنبال کرد.

    \item[تز ۳ (تکامل)]
    خمینی با خواندن پیش‌نویس متوجه شد ولایت فقیه در آن نیست — و این «نبود» باعث شد فکر کند ممکن است «اسلام» در ساختار نظام ضعیف بماند. تکامل فکری‌اش در این لحظه یک گام جلوتر رفت.

    \item[تز ۴ (فشار نخبگان)]
    بهشتی و حزب جمهوری اسلامی خمینی را متقاعد کردند. خمینی خودش ایدهٔ مجلس خبرگان را نداشت — اما وقتی بهشتی استدلال کرد، پذیرفت.
\end{description}
\end{principleBox}

% ────────────────────────────────────────
\section{مجلس خبرگان قانون اساسی: بازیگران و مذاکرات}
% ────────────────────────────────────────

\subsection{ترکیب مجلس}

\begin{figure}[H]
\centering
\begin{tikzpicture}
\pie[
    text=legend,
    radius=3.2,
    color={cPrimary!70, cGold!70, cGreen!60, cWarn!60, cGray!50},
    explode={0.05, 0.05, 0.05, 0.05, 0.05},
    sum=auto
]{
    55/روحانیون (عمدتاً حزب جمهوری اسلامی),
    8/حقوقدانان و تکنوکرات‌های مذهبی,
    4/نمایندگان اقلیت‌ها,
    3/نمایندگان چپ‌گرا یا ملی‌گرا,
    3/مستقل‌ها
}
\end{tikzpicture}
\caption{ترکیب مجلس خبرگان قانون اساسی (۷۳ عضو)}
\label{fig:experts-composition}
\end{figure}

\begin{warningBox}[اهمیت ترکیب]
با ۵۵ روحانی از ۷۳ عضو (۷۵\%)، نتیجهٔ مذاکرات از پیش قابل پیش‌بینی بود. منتقدان (از جمله بازرگان و مهندس سحابی) هشدار دادند اما کاری نتوانستند بکنند. آیت‌الله شریعتمداری نیز اعلام کرد ترکیب مجلس «نمایندهٔ واقعی ملت» نیست.\footnote{آیت‌الله شریعتمداری، مصاحبه با روزنامهٔ میزان، شهریور ۱۳۵۸.}
\end{warningBox}

\subsection{بازیگران کلیدی در مجلس خبرگان}

\begin{longtable}{%
    >{\raggedleft\arraybackslash}p{0.17\textwidth}
    >{\raggedleft\arraybackslash}p{0.18\textwidth}
    >{\raggedleft\arraybackslash}p{0.30\textwidth}
    >{\raggedleft\arraybackslash}p{0.25\textwidth}}
\toprule
\rowcolor{cPrimary!10}
\textbf{شخص} & \textbf{سمت} & \textbf{نقش در بحث ولایت فقیه} & \textbf{گرایش} \\
\midrule
\endhead

آیت‌الله منتظری &
رئیس مجلس &
حامی ولایت فقیه — اما با تعبیر «محدود» &
تندرو-اصولگرا (بعدها منتقد) \\
\midrule

آیت‌الله بهشتی &
نایب‌رئیس &
\keyword{معمار اصلی} ولایت فقیه در قانون اساسی &
تندرو-عملگرا \\
\midrule

عباسعلی عمید زنجانی &
عضو &
دفاع فقهی از ولایت فقیه &
تندرو \\
\midrule

حسن آیت &
عضو (غیرروحانی) &
مدافع سرسخت ولایت فقیه &
تندرو (بعدها ترور شد) \\
\midrule

مهندس سحابی &
عضو &
مخالف ولایت فقیه &
ملی-مذهبی \\
\midrule

رحمت‌الله مقدم مراغه‌ای &
عضو &
مخالف ولایت فقیه &
مستقل \\
\midrule

آیت‌الله طالقانی &
عضو (فوت مهر ۱۳۵۸) &
نگران تمرکز قدرت &
ملی-مذهبی \\

\bottomrule
\caption{بازیگران کلیدی مجلس خبرگان قانون اساسی}
\label{tab:experts-actors}
\end{longtable}

\subsection{مذاکرات: نبرد بر سر ولایت فقیه}

مذاکرات مجلس خبرگان از ۲۸ مرداد تا ۲۴ آبان ۱۳۵۸ طول کشید — حدود سه ماه. محور اصلی مناظرات، \keyword{اصل ولایت فقیه} بود.

\paragraph{موافقان ولایت فقیه — استدلال‌ها:}

\begin{enumerate}[label=\textcolor{cPrimary}{\bfseries\arabic*.}, rightmargin=2em]
    \item \textbf{استدلال فقهی:} ولایت فقیه امتداد ولایت ائمه (ع) است و بر اساس احادیث معتبر (مثل مقبولهٔ عمر‌بن‌حنظله) ثابت می‌شود.\footnote{صورت مشروح مذاکرات مجلس خبرگان، جلسهٔ ۲۵، ص ۴۴۸–۴۶۲.}
    
    \item \textbf{استدلال عملی:} بدون ولایت فقیه، «اسلامیت» نظام تضمین نمی‌شود. شورای نگهبان به‌تنهایی کافی نیست.
    
    \item \textbf{استدلال کاریزماتیک:} «امام» [خمینی] خود ولی فقیه است. نوشتن این اصل فقط «تثبیت واقعیت موجود» است.\footnote{بهشتی در مجلس خبرگان: «ما فقط آنچه هست را به زبان حقوقی درمی‌آوریم.» صورت مشروح مذاکرات، جلسهٔ ۲۷، ص ۵۱۰.}
\end{enumerate}

\paragraph{مخالفان ولایت فقیه — استدلال‌ها:}

\begin{enumerate}[label=\textcolor{cAccent}{\bfseries\arabic*.}, rightmargin=2em]
    \item \textbf{استدلال فقهی:} ولایت فقیه اجماع مراجع را ندارد. مراجع بزرگ (شریعتمداری، خویی، مرعشی نجفی) آن‌را نپذیرفته‌اند.
    
    \item \textbf{استدلال حقوقی:} ولایت فقیه با اصل حاکمیت ملی و جمهوریت تناقض دارد.
    
    \item \textbf{استدلال عملی:} تمرکز قدرت در دست یک فرد خطرناک است — حتی اگر آن فرد فقیه عادل باشد.\footnote{مقدم مراغه‌ای در مجلس خبرگان: «اگر ولی فقیه عادل نباشد چه؟ چه تضمینی هست؟» صورت مشروح مذاکرات، جلسهٔ ۲۸، ص ۵۳۵.}
    
    \item \textbf{استدلال تاریخی:} وعده‌های خمینی در پاریس با ولایت فقیه سازگار نیست.\footnote{سحابی در مجلس خبرگان: «مردم بر اساس وعده‌هایی رأی دادند... اگر از اول می‌گفتید ولایت فقیه، رأی مردم شاید متفاوت بود.» صورت مشروح مذاکرات، جلسهٔ ۳۰، ص ۵۸۲.}
\end{enumerate}

\begin{warningBox}[مرگ مشکوک طالقانی]
آیت‌الله سید محمود طالقانی، روحانی محبوب و امام جمعهٔ تهران، که به‌عنوان یکی از معدود صداهای منتقد تمرکز قدرت شناخته می‌شد، در ۱۹ شهریور ۱۳۵۸ — در اوج مذاکرات مجلس خبرگان — درگذشت. رسماً «سکتهٔ قلبی» اعلام شد، اما شایعات دربارهٔ \keyword{مرگ غیرطبیعی} هیچ‌گاه فروکش نکرد.\footnote{طالقانی هفته‌ها قبل از مرگ در سخنرانی‌ای گفته بود: «من خودم شخصاً به ولایت فقیه معتقدم اما نه آن‌طور که بعضی‌ها تفسیر می‌کنند... مردم باید آزاد باشند.» \en{Moin (1999), pp.\,218--219.}}

مرگ طالقانی عملاً مهم‌ترین صدای \keyword{اپوزیسیون درون‌نظامی} را خاموش کرد.
\end{warningBox}

\subsection{تصویب نهایی: آذر ۱۳۵۸}

قانون اساسی با ولایت فقیه در ۲۴ آبان ۱۳۵۸ به تصویب مجلس خبرگان رسید و در رفراندوم ۱۲ آذر ۱۳۵۸ با حدود ۹۹.۵\% آرا تأیید شد.\footnote{\en{Schirazi (1997), pp.\,36--40.} درصد بالای «آری» محل تردید منتقدان است — مشارکت حدود ۷۰\% بود اما فشار اجتماعی و فقدان آلترناتیو مؤثر بود.}

% ────────────────────────────────────────
\section{مقایسهٔ تفصیلی: پیش‌نویس در برابر مصوبه}
% ────────────────────────────────────────

\begin{longtable}{%
    >{\centering\arraybackslash}p{0.04\textwidth}
    >{\raggedleft\arraybackslash}p{0.15\textwidth}
    >{\raggedleft\arraybackslash}p{0.35\textwidth}
    >{\raggedleft\arraybackslash}p{0.35\textwidth}}
\toprule
\rowcolor{cPrimary!10}
\textbf{\#} & \textbf{موضوع} & \textbf{پیش‌نویس حبیبی-کاتوزیان} & \textbf{مصوبهٔ مجلس خبرگان} \\
\midrule
\endhead

۱ &
رهبری &
رئیس‌جمهور منتخب = بالاترین مقام &
ولی فقیه = بالاترین مقام؛ رئیس‌جمهور تابع \\
\midrule

۲ &
ولایت فقیه &
\keyword{وجود ندارد} &
اصل ۵ و اصل ۱۰۷–۱۱۲: ولایت امر و امامت امت \\
\midrule

۳ &
فرماندهی ارتش &
رئیس‌جمهور &
ولی فقیه \\
\midrule

۴ &
انتصاب رئیس قوهٔ قضاییه &
پیشنهاد مجلس + تأیید رئیس‌جمهور &
منصوب ولی فقیه \\
\midrule

۵ &
شورای نگهبان &
بررسی تطابق شرعی قوانین &
وتوی قوانین مجلس + نظارت بر انتخابات \\
\midrule

۶ &
آزادی احزاب &
بدون قید مذهبی &
مشروط به «موازین اسلامی» (اصل ۲۶) \\
\midrule

۷ &
آزادی مطبوعات &
تضمین‌شده &
مشروط به «عدم اخلال در مبانی اسلام» (اصل ۲۴) \\
\midrule

۸ &
حقوق زنان &
برابری کامل &
«با رعایت موازین اسلامی» (اصل ۲۰ و ۲۱) \\
\midrule

۹ &
حقوق اقلیت‌ها &
برابری شهروندی &
اقلیت‌های «رسمی» حقوق محدود — بهائیان حذف \\
\midrule

۱۰ &
بازنگری قانون اساسی &
مجلس مؤسسان + رفراندوم &
مجلس بازنگری + تأیید رهبر (اصل ۱۷۷) \\

\bottomrule
\caption{مقایسهٔ تفصیلی پیش‌نویس و مصوبهٔ نهایی قانون اساسی ۱۳۵۸}
\label{tab:constitution-detailed}
\end{longtable}

\begin{warningBox}[فاصلهٔ وعده تا عمل]
در \keyword{تمام} ده محور فوق، مصوبهٔ نهایی از پیش‌نویس (و بنابراین از وعده‌های پاریس) فاصله گرفته و به سمت تمرکز قدرت حرکت کرده است. این فاصله نه جزئی بلکه \keyword{ساختاری} است: دو سند دو نوع کاملاً متفاوت حکومت را تعریف می‌کنند.
\end{warningBox}

% ────────────────────────────────────────
\section{نقش خمینی در فرایند: «ناظر» یا «هدایت‌کننده»؟}
% ────────────────────────────────────────

\subsection{مداخلات مستقیم خمینی}

شواهدی وجود دارد که خمینی در مقاطع خاصی مستقیماً در مذاکرات مجلس خبرگان مداخله کرد:

\begin{enumerate}[label=\textcolor{cPrimary}{\bfseries\arabic*.}, rightmargin=2em]
    \item \textbf{پیام «عجله کنید»:} خمینی در شهریور ۱۳۵۸ پیامی به مجلس خبرگان فرستاد و خواست «کار را زودتر تمام کنید — دشمنان فرصت توطئه ندهید.» این فشار زمانی، فرصت بحث عمیق را محدود کرد.\footnote{صحیفهٔ امام، ج۹، ص ۲۲۰–۲۲۳.}
    
    \item \textbf{حمایت از اصل ولایت فقیه:} خمینی در سخنرانی‌هایی در شهریور و مهر ۱۳۵۸ صراحتاً از ولایت فقیه حمایت کرد — هرچند ادعا کرد «مجلس خبرگان خودشان تصمیم می‌گیرند.»\footnote{صحیفهٔ امام، ج۱۰، ص ۲۸–۳۵.}
    
    \item \textbf{سرکوب مخالفان:} هم‌زمان با مذاکرات مجلس خبرگان، حملات فیزیکی به دفاتر احزاب مخالف (حزب دموکرات کردستان، سازمان چریک‌های فدایی، حزب خلق مسلمان شریعتمداری) تشدید شد — که فضای ارعاب ایجاد کرد.\footnote{\en{Abrahamian (1999), pp.\,130--135.}}
\end{enumerate}

\subsection{شواهد «هدایت» غیرمستقیم}

\begin{itemize}[label=\textcolor{cConsolid}{$\blacktriangleright$}, rightmargin=2em]
    \item بهشتی مرتباً با دفتر خمینی (از طریق احمد خمینی) در ارتباط بود و پیش از طرح هر اصل مهم، «نظر امام» را جویا می‌شد.\footnote{رفسنجانی، \textit{عبور از بحران}، ص ۱۶۵–۱۸۰.}
    \item خمینی هیچ‌گاه علناً با اصل ولایت فقیه مخالفت نکرد — سکوت او خودش نوعی «تأیید» بود.
    \item اعضای مخالف مجلس خبرگان (مانند مقدم مراغه‌ای) بعداً شکایت کردند که «فضا طوری بود که مخالفت با ولایت فقیه مساوی بود با مخالفت با امام — و کسی جرأت این کار را نداشت.»\footnote{مصاحبهٔ مقدم مراغه‌ای با \en{BBC Persian}، ۱۳۸۰.}
\end{itemize}

% ────────────────────────────────────────
\section{واکنش‌ها: صداهای مخالف}
% ────────────────────────────────────────

\subsection{آیت‌الله شریعتمداری}

آیت‌الله سید کاظم شریعتمداری، از مراجع بزرگ تقلید و رقیب سنتی خمینی، صریحاً با ولایت فقیه مخالفت کرد:

\begin{noteBox}[مخالفت شریعتمداری]
\begin{itemize}[label=\textcolor{cGold}{$\blacktriangleright$}, rightmargin=2em]
    \item شریعتمداری گفت: «ولایت فقیه مبنای فقهی محکمی ندارد. اکثر فقها در طول تاریخ آن‌را نپذیرفته‌اند.»\footnote{مصاحبهٔ شریعتمداری با روزنامهٔ میزان، شهریور ۱۳۵۸.}
    \item او خواستار رفراندوم بر پیش‌نویس اولیه (بدون ولایت فقیه) شد.
    \item حزب خلق مسلمان (هوادار شریعتمداری) در تبریز تظاهرات برگزار کرد.
    \item نتیجه: سرکوب حزب خلق مسلمان، حصر خانگی شریعتمداری، و عزل او از مرجعیت (۱۳۶۱) — رویدادی بی‌سابقه در تاریخ شیعه.\footnote{\en{Bakhash (1984), pp.\,78--85;} همچنین: \en{Moin (1999), pp.\,222--228.}}
\end{itemize}
\end{noteBox}

\subsection{بازرگان و نهضت آزادی}

بازرگان ناتوان از جلوگیری از تصویب ولایت فقیه بود. او در خاطراتش نوشت:

\begin{warningBox}[روایت بازرگان]
\textit{«قانون اساسی‌ای که تصویب شد، قانون اساسی‌ای نیست که ما برایش انقلاب کردیم. ما جمهوری می‌خواستیم، ولایت [فقیه] گرفتیم.»}\footnote{مهدی بازرگان، \textit{انقلاب ایران در دو حرکت}، ص ۱۸۵.}
\end{warningBox}

\subsection{روشنفکران و فعالان سکولار}

\begin{itemize}[label=\textcolor{cAccent}{$\blacktriangleright$}, rightmargin=2em]
    \item \textbf{هدایت‌الله متین‌دفتری} (نوهٔ مصدق، رهبر جبههٔ دموکراتیک ملی): ولایت فقیه را «بازگشت به استبداد» خواند.\footnote{\en{Bakhash (1984), p.\,80.}}
    \item \textbf{کریم سنجابی} (رهبر جبههٔ ملی): «این قانون اساسی روح مشروطه‌خواهی ایران را نقض می‌کند.»\footnote{کریم سنجابی، \textit{خاطرات}، (منتشرشده پس از فوت).}
    \item \textbf{دکتر هما کاتوزیان}: «پیش‌نویسی که ما نوشتیم، بازتاب وعده‌هایی بود که خمینی در پاریس داد. آنچه تصویب شد، نقض آن وعده‌هاست.»\footnote{\en{Katouzian, ``The Short-Lived Hopes,'' various conference papers.}}
\end{itemize}

% ────────────────────────────────────────
\section{گاه‌شمار تطبیقی: از پیش‌نویس تا تصویب}
% ────────────────────────────────────────

\begin{figure}[H]
\centering
\begin{tikzpicture}[
    x=1.3cm,
    event/.style={
        rectangle,
        rounded corners=2pt,
        draw=#1!70,
        fill=#1!8,
        text width=3.4cm,
        font=\scriptsize,
        align=center,
        inner sep=3pt
    },
    date/.style={
        font=\scriptsize\bfseries,
        color=cPrimary
    },
    line/.style={
        cPrimary!40,
        line width=2pt
    }
]

\draw[line] (0,0) -- (15,0);

\node[event=cGreen, above=0.7cm] (e1) at (0,0) {پیش‌نویس حبیبی-کاتوزیان\\منتشر می‌شود};
\draw[-{Stealth[length=2.5pt]}, cGreen!50] (0,0) -- (e1);
\node[date] at (0,0.25) {۲۶ خرداد ۵۸};
\fill[cGreen] (0,0) circle (2.5pt);

\node[event=cGold, below=0.7cm] (e2) at (3,0) {خمینی: «خوب است»\\پیشنهاد رفراندوم};
\draw[-{Stealth[length=2.5pt]}, cGold!50] (3,0) -- (e2);
\node[date] at (3,0.25) {تیر ۵۸};
\fill[cGold] (3,0) circle (2.5pt);

\node[event=cWarn, above=0.7cm] (e3) at (5.5,0) {حزب جمهوری اسلامی:\\«مجلس خبرگان لازم است»};
\draw[-{Stealth[length=2.5pt]}, cWarn!50] (5.5,0) -- (e3);
\node[date] at (5.5,0.25) {تیر–مرداد ۵۸};
\fill[cWarn] (5.5,0) circle (2.5pt);

\node[event=cAccent, below=0.7cm] (e4) at (8,0) {انتخابات مجلس خبرگان\\پیروزی روحانیون (۷۵\%)};
\draw[-{Stealth[length=2.5pt]}, cAccent!50] (8,0) -- (e4);
\node[date] at (8,0.25) {۱۲ مرداد ۵۸};
\fill[cAccent] (8,0) circle (2.5pt);

\node[event=cPrimary, above=0.7cm] (e5) at (10.5,0) {مذاکرات مجلس خبرگان\\تصویب ولایت فقیه};
\draw[-{Stealth[length=2.5pt]}, cPrimary!50] (10.5,0) -- (e5);
\node[date] at (10.5,0.25) {مرداد–آبان ۵۸};
\fill[cPrimary] (10.5,0) circle (2.5pt);

\node[event=cAccent, below=0.7cm] (e6) at (13.5,0) {رفراندوم قانون اساسی\\۹۹.۵\% آری};
\draw[-{Stealth[length=2.5pt]}, cAccent!50] (13.5,0) -- (e6);
\node[date] at (13.5,0.25) {۱۲ آذر ۵۸};
\fill[cAccent] (13.5,0) circle (2.5pt);

\end{tikzpicture}
\caption{گاه‌شمار: از پیش‌نویس تا تصویب قانون اساسی (خرداد–آذر ۱۳۵۸)}
\label{fig:constitution-timeline}
\end{figure}

% ────────────────────────────────────────
\section{جمع‌بندی فصل هشتم}
% ────────────────────────────────────────

\begin{principleBox}[خلاصهٔ فصل ۸]
\begin{itemize}[label=\textcolor{cPrimary}{$\blacksquare$}, rightmargin=2em]
    \item پیش‌نویس حبیبی-کاتوزیان یک جمهوری دموکراتیک بدون ولایت فقیه بود — کاملاً سازگار با وعده‌های پاریس.
    \item خمینی ابتدا پیش‌نویس را تأیید کرد اما سپس مسیر را به سمت مجلس خبرگان تغییر داد.
    \item ترکیب مجلس خبرگان (۷۵\% روحانی) نتیجه را از پیش تعیین کرد.
    \item بهشتی «معمار اصلی» ولایت فقیه در قانون اساسی بود — اما خمینی فعالانه حمایت کرد.
    \item صداهای مخالف (شریعتمداری، بازرگان، سحابی) سرکوب یا حاشیه‌نشین شدند.
    \item \keyword{ارزیابی:} این بزنگاه بیش از هر رویداد دیگری با \keyword{تز دوم (خدعه)} و \keyword{تز چهارم (فشار نخبگان)} سازگار است. تز اول (صداقت + اجبار) در توضیح این فرایند \keyword{ضعیف‌ترین} عملکرد را دارد — چون هیچ «بحران اضطراری» در خرداد تا آذر ۱۳۵۸ وجود نداشت که تغییر مسیر را «اجبار» کند.
\end{itemize}
\end{principleBox}


% ============================================================
%  فصل ۹: ترورهای گروه فرقان و بحران امنیتی
% ============================================================
\chapter{ترورهای گروه فرقان و بحران امنیتی}

\begin{infoBox}[چکیدهٔ فصل]
گروه فرقان، یک سازمان مسلح کوچک با ایدئولوژی التقاطی (ترکیب اسلام و مارکسیسم)، بین اردیبهشت تا آذر ۱۳۵۸ چند ترور سیاسی انجام داد که مهم‌ترین آن‌ها قتل آیت‌الله مرتضی مطهری، سرتیپ قرنی، و آیت‌الله مفتح بود. این ترورها نقشی کلیدی در \keyword{فضاسازی امنیتی} ایفا کردند و به بهانهٔ مقابله با «ضدانقلاب»، فرایند تمرکز قدرت و محدودیت آزادی‌ها تسریع شد. این فصل به بررسی ماهیت فرقان، ترورها، و پیامدهای سیاسی آن‌ها می‌پردازد.
\end{infoBox}

% ────────────────────────────────────────
\section{گروه فرقان: ماهیت و ایدئولوژی}
% ────────────────────────────────────────

\subsection{تأسیس و رهبری}

گروه فرقان در اواخر دهه ۱۳۵۰ توسط \keyword{اکبر گودرزی} تأسیس شد. گودرزی دانشجوی مهندسی و تحت تأثیر اندیشه‌های علی شریعتی بود اما خوانش رادیکال‌تری داشت:\footnote{\en{Abrahamian (1989),} \textit{Radical Islam: The Iranian Mojahedin}, I.\,B.\,Tauris, pp.\,182--190. همچنین: \en{Bakhash (1984), pp.\,62--65.}}

\begin{itemize}[label=\textcolor{cAccent}{$\blacktriangleright$}, rightmargin=2em]
    \item فرقان «روحانیت» را مانع تحقق «اسلام واقعی» می‌دانست.
    \item معتقد بود خمینی و روحانیون دارند انقلاب را «مصادره» می‌کنند.
    \item خواستار یک «حکومت اسلامی بدون روحانیون» بود — «اسلام منهای آخوند».
    \item تعداد اعضا بسیار محدود بود — شاید ۵۰ تا ۱۰۰ نفر.
\end{itemize}

\begin{noteBox}[ایدئولوژی التقاطی فرقان]
فرقان از نظر فکری التقاطی بود: ترکیبی از شریعتی (اسلام انقلابی)، مارکسیسم (مبارزهٔ طبقاتی)، و ضدروحانیت (نقد نهاد مرجعیت). این التقاط باعث شد فرقان هم از طرف روحانیون (به‌خاطر ضدروحانیت) و هم از طرف چپ‌های مارکسیست (به‌خاطر اسلام‌گرایی) طرد شود — و عملاً \keyword{هیچ پایگاه اجتماعی} نداشت.
\end{noteBox}

% ────────────────────────────────────────
\section{ترورها: گاه‌شمار}
% ────────────────────────────────────────

\begin{longtable}{%
    >{\centering\arraybackslash}p{0.05\textwidth}
    >{\raggedleft\arraybackslash}p{0.14\textwidth}
    >{\raggedleft\arraybackslash}p{0.20\textwidth}
    >{\raggedleft\arraybackslash}p{0.22\textwidth}
    >{\raggedleft\arraybackslash}p{0.29\textwidth}}
\toprule
\rowcolor{cAccent!10}
\textbf{\#} & \textbf{تاریخ} & \textbf{قربانی} & \textbf{سمت} & \textbf{پیامد سیاسی} \\
\midrule
\endhead

۱ &
۲۳ فروردین ۱۳۵۸ &
سرتیپ محمدولی قرنی &
رئیس ستاد ارتش &
تشدید پاکسازی ارتش، تقویت سپاه \\
\midrule

۲ &
۱۲ اردیبهشت ۱۳۵۸ &
آیت‌الله مرتضی مطهری &
عضو شورای انقلاب، اندیشمند &
شوک ملی، تقویت روایت «توطئهٔ ضدانقلاب» \\
\midrule

۳ &
۲۹ آذر ۱۳۵۸ &
آیت‌الله محمد مفتح &
استاد دانشگاه، عضو شورای انقلاب &
تشدید سرکوب گروه‌های مسلح \\

\bottomrule
\caption{ترورهای اصلی گروه فرقان (۱۳۵۸)}
\label{tab:forqan-assassinations}
\end{longtable}

\subsection{ترور مطهری: شوک و پیامدها}

ترور آیت‌الله مرتضی مطهری در شب ۱۲ اردیبهشت ۱۳۵۸ (یک روز پس از روز معلم) مهم‌ترین ترور فرقان بود. مطهری از نزدیک‌ترین شاگردان خمینی، عضو شورای انقلاب، و یکی از محبوب‌ترین روحانیون-روشنفکر بود.\footnote{\en{Moin (1999), pp.\,214--216.}}

\begin{warningBox}[واکنش خمینی به ترور مطهری]
خمینی در واکنش به ترور مطهری گفت:

\textit{«فرزند عزیزم را کشتند... این‌ها دشمنان اسلامند... دست‌های پنهان در کارند.»}\footnote{صحیفهٔ امام، ج۷، ص ۱۱۰–۱۱۵.}

و سپس:

\textit{«همهٔ گروه‌هایی که اسلحه دارند، باید تحویل دهند... کسی حق ندارد اسلحه داشته باشد مگر سپاه و ارتش و کمیته‌ها.»}\footnote{همان، ص ۱۲۰.}

مدافعان تز دوم تأکید می‌کنند که خمینی از ترور مطهری برای \keyword{خلع سلاح عمومی} استفاده کرد — اقدامی که بیشتر علیه مجاهدین خلق و چریک‌های فدایی بود تا فرقان.
\end{warningBox}

% ────────────────────────────────────────
\section{پیامدهای سیاسی ترورها}
% ────────────────────────────────────────

\subsection{تقویت نهادهای انقلابی}

ترورهای فرقان بهانهٔ مهمی برای تقویت نهادهای امنیتی-انقلابی فراهم کرد:

\begin{enumerate}[label=\textcolor{cPrimary}{\bfseries\arabic*.}, rightmargin=2em]
    \item \textbf{تقویت سپاه پاسداران:} سپاه اعلام کرد فقط او توانایی مقابله با «ضدانقلاب» را دارد. بودجه و نیروی سپاه افزایش یافت.\footnote{\en{Bakhash (1984), pp.\,63--65.}}
    
    \item \textbf{تقویت دادگاه‌های انقلاب:} خلخالی از ترورها به‌عنوان توجیه سرعت‌بخشی به محاکمات استفاده کرد.
    
    \item \textbf{محدودسازی مطبوعات:} روزنامه‌هایی که از «آزادی گروه‌ها» دفاع می‌کردند، تحت فشار قرار گرفتند. روزنامهٔ \textit{آیندگان} در مرداد ۱۳۵۸ توقیف شد.\footnote{\en{Abrahamian (1999), pp.\,132--135.}}
    
    \item \textbf{خلع سلاح گروه‌ها:} فرمان خمینی برای تحویل اسلحه، ابزاری برای ضعیف‌سازی تمام گروه‌های مسلح غیرحکومتی شد.
\end{enumerate}

\subsection{روایت «توطئهٔ گسترده»}

حلقهٔ درونی خمینی از ترورهای فرقان برای ساختن روایت \keyword{توطئهٔ گسترده} بهره برد:

\begin{itemize}[label=\textcolor{cConsolid}{$\blacktriangleright$}, rightmargin=2em]
    \item ترورها به‌عنوان بخشی از یک «طرح آمریکایی-صهیونیستی» معرفی شدند.
    \item فرقان با ساواک، \en{CIA}، و گروه‌های مارکسیست «مرتبط» دانسته شد — هرچند شواهد قاطعی ارائه نشد.\footnote{اعترافات تلویزیونی اعضای فرقان (قبل از اعدام) با شکنجه به‌دست آمده بود و اعتبار آن‌ها محل تردید است. \en{Abrahamian (1999), pp.\,135--140.}}
    \item هر گروه مخالف — حتی لیبرال‌ها — در معرض اتهام «همکاری با ضدانقلاب» قرار گرفت.
\end{itemize}

\begin{warningBox}[«ضدانقلاب» — مفهومی کشدار]
مفهوم \keyword{«ضدانقلاب»} (\rl{ضد انقلاب}) پس از ترورهای فرقان به اصلی‌ترین ابزار تبلیغاتی حلقهٔ درونی تبدیل شد. این مفهوم عمداً \keyword{مبهم} نگاه داشته شد تا بتواند شامل هر مخالفی بشود: از فرقان و مجاهدین گرفته تا بازرگان و شریعتمداری. هر کسی که با تمرکز قدرت مخالفت می‌کرد، بالقوه «ضدانقلاب» بود.
\end{warningBox}

% ────────────────────────────────────────
\section{تفسیرهای رقیب از ترورهای فرقان}
% ────────────────────────────────────────

\begin{principleBox}[چهار خوانش از نقش ترورها]
\begin{description}[style=nextline, font=\bfseries\color{cPrimary}]
    \item[تز ۱ (صداقت + اجبار)]
    ترورها تهدید واقعی بودند. خمینی مجبور شد امنیت را بر آزادی مقدم بدارد. تقویت سپاه و محدودیت‌ها \keyword{واکنش طبیعی} به تهدید بود.

    \item[تز ۲ (خدعه)]
    ترورها بهانه‌ای «طلایی» بودند. حلقهٔ درونی از آن‌ها برای پیشبرد دستور کار از پیش طراحی‌شدهٔ تمرکز قدرت استفاده کرد. ممکن است حتی برخی اقدامات تحریک‌آمیز عمدی بوده باشند.\footnote{این ادعا (تحریک عمدی) شاهد مستقیمی ندارد و در حد حدس باقی می‌ماند.}

    \item[تز ۳ (تکامل)]
    ترورها خمینی را متقاعد